% ============================================================================
% Część 6 przewodnika: wersja webowa (rozmowa w przeglądarce).
% Wczytywana przez \input z przewodnik.pl.tex po czesc-5-cwiczenia; zawiera
% wyłącznie treść od \section.
%
% Rodzaj dokumentu: przewodnik, nie artykuł naukowy w rozumieniu standardu
% (docs/research/00-RESEARCH-DOCUMENTATION.md, „Documents that borrow the house
% style"). Nie wprowadza żadnego faktu spoza źródeł. Źródła, z których czerpie:
% docs/architecture/web-app-plan.md (kontrakt §10 i decyzje §9), ADR-0076,
% ADR-0077, ADR-0078, docs/runbooks/web-app-cost-controls.md, docs/legal/
% (szkice i CHECKLIST), docs/ux/UI-UX.md, protokół
% src/ExitInterviewAgent.Agent/Protocol/interview-protocol.pl.v1.json oraz kod
% w src/ExitInterviewAgent.InterviewService i web/app. Gdzie kod i szkic
% prawny się rozjeżdżają, rozdział to mówi wprost.
%
% Stan, który rozdział opisuje: wersja webowa jest ZAIMPLEMENTOWANA, ale NIE
% wdrożona, bez przeglądu prawnego i bez przebiegu z prawdziwym modelem (W9).
% Wszystko, co wymaga tych rzeczy, jest oznaczone jako NOT RUN.
%
% Ćwiczenia w części 4 są oparte o komendy uruchomione w sesji pisania
% (dotnet SDK 10.x z pakietu apt, model testowy „mock”, bez sieci i bez kluczy).
% Każde wyjście w ramce Verbatim pochodzi z tej sesji; skrócone tylko tam,
% gdzie widać „...”.
% ============================================================================

{\sloppy\emergencystretch=4em
\section{Wersja webowa: rozmowa w przeglądarce}
\label{sec:webapp}

Wersja webowa przenosi wywiad z terminala do przeglądarki. Dla osoby, która
nie chce używać terminala, to jedyna droga do rozmowy. Logika rozmowy jest ta
sama co w CLI: maszyna stanów z protokołu i ten sam prowadzący, a warstwa
webowa dodaje sesje, konta, kredyty, płatność i limity kosztów. Nie dodaje
nowych zachowań rozmowy.

\finding{Stan tej części, przeczytaj najpierw.} Wersja webowa jest
\emph{zaimplementowana}, ale \emph{nie wdrożona}. Nie ma przeglądu prawnego
(lista pytań dla prawnika ma status NOT REVIEWED przy każdej pozycji). Nie
przeprowadzono przebiegu z prawdziwym modelem, więc nie wiadomo, jak brzmi
rozmowa po polsku z prawdziwym dostawcą, ile kosztuje jedna rozmowa ani jaka
ma być cena. Ten rozdział nie podaje tych liczb. Nie wolno przyjmować
płatności ani prowadzić rozmowy z prawdziwymi ludźmi przed tymi krokami.

\subsection{Stan poszczególnych zadań}

Plan (\repopath{docs/architecture/web-app-plan.md}, §7) dzieli budowę na zadania
W1--W10. Poniżej stan według repozytorium, a nie według planu.

\begin{longtable}{@{}p{0.1\linewidth}p{0.52\linewidth}p{0.32\linewidth}@{}}
\toprule
Zadanie & Co zawiera & Stan \\
\midrule
\endhead
W1 & ADR-y decyzji webowych (ADR-0076, ADR-0077, ADR-0078) i decyzje w planie &
     w repozytorium \\
W2 & Sesje w \texttt{InterviewService} (\texttt{/api/v1/interviews}), pamięć
     procesu, wygaszanie po 30 min, testy &
     zaimplementowane \\
W3 & Księga kredytów, zdarzenia płatności, webhook z podpisem, zwrot kredytu,
     testy na dostawcy testowym &
     zaimplementowane; \textbf{test z prawdziwym Stripe: NOT RUN} \\
W4 & Wyłącznik, limity na konto i adres, bramka e-mail, dzienny limit, metryki
     bez etykiet &
     zaimplementowane; \textbf{bramka e-mail blokuje wszystkie starty} (zob.
     niżej) \\
W5, W6 & Ścieżka BFF i strona \texttt{/interview} (zgoda, czat, wynik, usunięcie),
     po polsku i po angielsku &
     w repozytorium \\
W7 & Test przeglądarkowy całej drogi na dostawcy testowym &
     poza zakresem tego rozdziału (zob. \repopath{tests/e2e}) \\
W8 & Szkice polityki prywatności, regulaminu, tekstu zgody i lista pytań dla
     prawnika &
     \textbf{szkice; przegląd prawny: NOT RUN} \\
W9 & Przebieg z prawdziwym modelem: 10 rozmów po polsku, koszt, jakość, cena &
     \textbf{NOT RUN} (wymaga klucza i pieniędzy właściciela) \\
W10 & Wdrożenie na Fly (najpierw staging), limit wydatków w konsoli dostawcy,
      runbook &
     \textbf{NOT RUN} (zakaz w AGENTS.md: nic nie jest wdrożone) \\
\bottomrule
\end{longtable}

\subsection{Dla byłego pracownika: jak wygląda droga}

Droga składa się z pięciu kroków: konto, kredyt, zgoda, rozmowa i wynik.
Każdy krok opisuje to, co widać na ekranie, zgodnie z tabelą ekranów w
\repopath{docs/ux/UI-UX.md}.

\subsubsection{Konto: rejestracja i logowanie}

Konto zakłada się w usłudze tożsamości (authservice) na stronie
\texttt{/register}. Po rejestracji przychodzi e-mail z linkiem do
\texttt{/verify-email}. Logowanie działa na \texttt{/login}; jeśli konto ma
dwuetapowe logowanie, drugi krok przyjmuje kod z aplikacji albo jeden kod
odzyskiwania. Strona nie widzi tokenu dostępu: trzyma go BFF w ciasteczku sesji.

\uwaga{Bramka e-mail. Konfiguracja domyślnie wymaga zweryfikowanego e-maila
(\texttt{Interviews:RequireVerifiedEmail=true}), a authservice nie wpisuje
jeszcze \texttt{email\_\allowbreak{}verified} do tokenu (ADR-0078). Skutek: przy tej
konfiguracji każdy start rozmowy kończy się odmową \texttt{403
email\_\allowbreak{}not\_\allowbreak{}verified}. Ta odmowa jest zamierzona: zamknięte ustawienie domyślne.
Wyłączenie bramki na produkcji „żeby starty działały” jest błędem. Poprawką jest
zmiana w authservice.}

\subsubsection{Kredyt: ile rozmów i jak je kupić}

Rozmowa zużywa jeden kredyt. Ekran startu pokazuje wybór języka (pierwszy raz
język przeglądarki, potem wybór osoby), pytanie o czas w firmie (sześć pasm:
do 6 miesięcy, 6 miesięcy--1 rok, 1--3 lata, 3--5 lat, 5--10 lat, ponad 10 lat),
liczbę posiadanych kredytów i przycisk \emph{Dalej} (wymaga jednego kredytu).
Przycisk \emph{Kup jedną rozmowę} przenosi na stronę dostawcy płatności
(\texttt{POST /api/v1/checkout}, ilość 1--10). Kredyt pojawia się dopiero po
zdarzeniu płatności od dostawcy, które przychodzi na webhook i jest sprawdzone
podpisem. Strona sama nie dodaje kredytu.

\uwaga{W tej konfiguracji płatność jest wyłączona. Domyślnie
\texttt{Billing:Provider} ma wartość \texttt{none}: płatność i webhook odpowiadają
503 \texttt{billing\_\allowbreak{}disabled}. Cena nie ma wartości domyślnej
(\texttt{Billing:PriceMinorUnits}); bez niej checkout też odmawia. Cena ma być
ustawiona po przebiegu W9.}

\subsubsection{Zgoda}

Przed pierwszym pytaniem pojawia się pełny tekst zgody, w języku rozmowy. Mówi,
że rozmowę prowadzi AI, że wszystko, co wpiszesz, idzie do dostawcy modelu
(Anthropic), że rozmowa jest trzymana w pamięci serwera i kasowana po zakończeniu,
że wynik jest czytelny przez 30 minut, że możesz nacisnąć \emph{Zatrzymaj i usuń}
w każdej chwili, że szkice nie są faktami ani poradą prawną i że usługa nie
sprawdza zatrudnienia. Dopiero po zaznaczeniu pola i \emph{Zacznij rozmowę}
zużywa się kredyt. \emph{Wstecz} nie zużywa niczego i niczego nie zapisuje.

\finding{Rozbieżność, którą trzeba zamknąć przed startem.} Szkic zgody
(\repopath{docs/legal/consent-wording.md}) przewiduje trzy pola: A (AI i odbiorca,
wymagane), B (regulamin i polityka, wymagane) i C (opcjonalne zgłoszenie do
sygnałów). W kodzie strony jest \emph{jedno} pole. Pola C w webowej rozmowie
nie ma, bo ścieżka zgłoszenia z webu nie jest zbudowana (kod mówi o niej
„later”). Szkic prawny i kod trzeba zgodzić, zanim pojawi się prawdziwy
użytkownik.

\subsubsection{Rozmowa: co robi prowadzący}

Rozmowa ma sześć tematów, w tej kolejności: \emph{Wdrożenie}, \emph{Zarządzanie},
\emph{Rozwój}, \emph{Wynagrodzenie a obietnice}, \emph{Kultura}, \emph{Powód
odejścia}. Prowadzący zadaje po jednym pytaniu z tematu, a każda odpowiedź
przechodzi przez tę samą kolejność reguł (protokół, pliki \texttt{interview-protocol*.json}).
Reguły wybierają następny krok, a model tylko formułuje pytania tematów.
Najważniejsze zachowania:

\begin{itemize}[leftmargin=*,itemsep=2pt]
  \item \textbf{Pierwsza wiadomość} mówi, że prowadzący jest AI, czym jest zapis,
        że można przerwać w każdej chwili, i prosi o zgodę. Niejasna odpowiedź
        na tę prośbę dostaje jedno ponowne pytanie; druga niejasna odpowiedź
        kończy rozmowę bez zapisu.
  \item \textbf{Krótka odpowiedź z treścią} (na przykład „słabe”) dostaje jedno
        dopytanie: „Jaka jedna konkretna sytuacja albo moment najlepiej to
        ilustruje?”. Krótka odpowiedź nie jest zamykana.
  \item \textbf{Trzy odpowiedzi z rzędu bez treści} (do trzech słów i bez
        wyraźnej oceny) kończą rozmowę z zapisem na podstawie tego, co już padło.
  \item \textbf{Imię osoby} w odpowiedzi nie jest zapisywane: prowadzący raz
        na temat prosi o opis przez rolę (przełożony, współpracownik) i
        podaje, że imion nie zapisuje.
  \item \textbf{Wrogość}: prowadzący uznaje frustrację i nie naciska (bez
        dopytania i bez przekierowania). Druga taka odpowiedź w całej rozmowie
        kończy ją z zapisem.
  \item \textbf{Poważna relacja} (przykład: zwolnienie, mobbing) otwiera
        pogłębienie. Menu ma pięć kierunków w stałej kolejności: co się
        wydarzyło, kiedy i jak często, kto i w jakiej roli, co ta osoba
        zrobiła i jaka była reakcja, jak się to skończyło i jakie to ma
        znaczenie. Na jeden temat są maksymalnie cztery pogłębienia. Kierunek,
        który odpowiedź już opisała, jest pomijany.
  \item \textbf{Sprzeczność} w odpowiedziach dostaje jedno doprecyzowanie na temat.
  \item \textbf{Prośba o zakończenie} (na przykład „możemy już skończyć?”) to
        nie wycofanie zgody. Prowadzący kończy grzecznie i \emph{zapisuje}
        rozmowę. Nie zadaje kolejnego pytania.
  \item \textbf{Wycofanie zgody} w trakcie rozmowy kończy ją bez zapisu: nic
        z niej nie zostaje, rekord się nie tworzy, a kredyt \emph{nie wraca}
        (kontrakt §10: status \texttt{stopped} nie zwraca kredytu).
  \item \textbf{Limity rozmowy}: najwyżej 60 tur prowadzącego, 90 wywołań modelu
        i 90 000 szacowanych tokenów (protokół). Po ich wyczerpaniu rozmowa
        kończy się zapisem.
\end{itemize}

Każda tura to jedna odpowiedź w osobnym żądaniu. Strona pokazuje wskaźnik
„prowadzący pisze…”, aż przyjdzie następna tura. Długość odpowiedzi jest
ograniczona do 2000 znaków; dłuższa jest odrzucana z \texttt{422}, a tekst
zostaje w polu.

\subsubsection{Zakończenie i \emph{Zatrzymaj i usuń}}

Przycisk \emph{Zatrzymaj i usuń} jest w trakcie rozmowy i w wyniku. Pyta raz, a
potem wywołuje \texttt{DELETE /api/v1/interviews/\{id\}}, który kasuje rozmowę
i wynik z pamięci. Kolejny odczyt tej sesji daje 404. Usunięcie \emph{nie}
zwraca kredytu: wybór był Twój. To samo robi \emph{Usuń wszystko teraz} na
ekranie wyniku. Rozmowa, która wygasła po 30 minutach bez aktywności, daje
410 \texttt{gone} przez kolejne 30 minut. Wygaśnięcie bezczynności nie zwraca
kredytu w obecnym kodzie, bo nie jest błędem usługi (zob. „Ograniczenia” w
części dla operatora).

\subsubsection{Wynik: rekord i szkice}

Po zakończeniu rozmowy strona pokazuje:

\begin{itemize}[leftmargin=*,itemsep=2pt]
  \item na górze: zdanie prawne i ostrzeżenie, że to szkice;
  \item \textbf{rekord} czytany jako zdania, po jednym bloku na temat: ocena,
        pewność, krótkie cytaty z Twoich słów (imiona zamaskowane). Surowy JSON
        jest tylko do pobrania;
  \item \textbf{kafelki} (szkice tekstów do publikacji), każdy z przyciskiem
        \emph{Kopiuj}: \emph{Glassdoor}, \emph{Opinia Google}, \emph{Reddit},
        krótka notatka, przegląd i fakty. Kopiowanie wkłada na schowek sam tekst,
        bez etykiety i bez zdania prawnego. Gdy schowek jest niedostępny, strona
        mówi to wprost;
  \item pobieranie: rekord (JSON), kafelki jako \texttt{tiles.html} (budowany
        w przeglądarce, bez skryptu i bez zasobów sieciowych) i kafelki (JSON);
  \item \emph{Usuń wszystko teraz}.
\end{itemize}

Wynik zostaje w pamięci serwera 30 minut po zakończeniu rozmowy, potem znika.
Strona nie oferuje zachowania go nigdzie indziej. Kafelki są szkicami: nikt ich
nie publikuje, usługa nie wysyła ich dalej, a publikujesz je sam, na własną
odpowiedzialność (szkic regulaminu, \repopath{docs/legal/terms.pl.md}).

\subsubsection{Gdy coś się nie uda}

Komunikaty błędów mapuje jedna funkcja (\texttt{failureFor} w
\repopath{web/app/lib/interview-api.ts}). W skrócie:

\begin{longtable}{@{}p{0.2\linewidth}p{0.5\linewidth}p{0.22\linewidth}@{}}
\toprule
Odpowiedź & Co widzisz & Co dalej \\
\midrule
\endhead
402 \texttt{payment\_\allowbreak{}required} & brak kredytu; kup jeden & ekran startu \\
409 \texttt{interview\_\allowbreak{}in\_\allowbreak{}progress} & masz otwartą rozmowę; dokończ albo usuń & zgoda \\
410 \texttt{gone} & rozmowa wygasła po 30 min bez aktywności; kredyt nie wraca & start \\
404 \texttt{not\_found} & po restarcie usługi rozmowa nie istnieje; kredyt nie wraca sam (zob. wyżej) & start \\
422 \texttt{reply\_\allowbreak{}invalid} & pusta odpowiedź albo ponad 2000 znaków; tekst zostaje & czat \\
429 \texttt{rate\_\allowbreak{}limited} & za dużo żądań; czekaj, strona nie ponawia sama & tam, gdzie byłeś \\
503 \texttt{provider\_\allowbreak{}unavailable} & model nie odpowiada; rozmowa trwa, tekst zostaje & czat \\
503 \texttt{interviews\_\allowbreak{}disabled} & wywiady wstrzymane; kredyt nie jest zużywany & tam, gdzie byłeś \\
403 \texttt{email\_\allowbreak{}not\_\allowbreak{}verified} & potwierdź adres e-mail (zob. wyżej) & start \\
\bottomrule
\end{longtable}

\subsection{Co dzieje się z danymi}

Ta część odpowiada na pytanie, gdzie są dane, i opisuje to, co kod robi, a
nie to, co planowano. Każda deklaracja w szkicu polityki
(\repopath{docs/legal/privacy-policy.pl.md}, §3) jest tu zestawiona z tym, co
faktycznie działa.

\begin{longtable}{@{}p{0.22\linewidth}p{0.36\linewidth}p{0.36\linewidth}@{}}
\toprule
Dane & Gdzie są & Jak długo \\
\midrule
\endhead
Tekst rozmowy & tylko w pamięci procesu serwera; wysyłany do dostawcy modelu
(Anthropic) przy każdej turze & do końca rozmowy albo 30 min bez aktywności \\
Wynik (rekord, kafelki) & tylko w pamięci procesu; zwracany osobie & 30 min po
zakończeniu \\
Sesja (identyfikator, właściciel, status, znaczniki czasu, liczniki tokenów) &
baza \texttt{interviewdb} & nie ustalono; w szkicu polityki pole ma wartość
\texttt{[okres, do ustalenia]} \\
Kredyty & księga kredytów w bazie: wiersze dopisywane, nic nie jest
nadpisywane; saldo to suma wierszy & według decyzji prawnej (L5, L6) \\
Zdarzenia płatności & baza, jeden wiersz na identyfikator zdarzenia & jw. \\
Logi i metryki & bez tekstu rozmowy, bez adresu e-mail, bez identyfikatora
konta; same liczby & jw. \\
\bottomrule
\end{longtable}

Trzy rzeczy, które warto wiedzieć wprost:

\begin{itemize}[leftmargin=*,itemsep=2pt]
  \item \textbf{Tekst nie trafia do bazy, bo sesje nie mają tam tabeli.} Sesje
        są w \texttt{InterviewSessionStore} w pamięci procesu. Test
        \texttt{No\_\allowbreak{}interview\_\allowbreak{}text\_\allowbreak{}reaches\_\allowbreak{}the\_\allowbreak{}logs\_\allowbreak{}or\_\allowbreak{}the\_\allowbreak{}spans} w
        \texttt{InterviewEndpointTests} wstawia kanarka do odpowiedzi i sprawdza,
        że nie ma go w logach ani w śladach. Nie ma osobnego testu, który
        przeszukuje bazę po całej rozmowie, bo tekst tam nie jest zapisywany.
  \item \textbf{Rozmowa jest u dostawcy modelu.} Usługa wysyła ją do dostawcy
        modelu (Anthropic) na kluczu serwera. Jak długo dostawca ją trzyma i czy
        jej używa, zależy od jego warunków dla konta API. Projekt tych warunków
        nie zweryfikował (L5 w \repopath{docs/legal/CHECKLIST.md}).
  \item \textbf{Zgłoszenie do sygnałów pracodawców} z rozmowy webowej w tej
        wersji nie istnieje. Ścieżka CLI i MCP (opt-in) istnieje w innym miejscu
        repozytorium. Gdy webowe zgłoszenie powstanie, będzie wymagało zgody C
        i testu, że bez zgody nie trafia nigdzie (CHECKLIST, sekcja C).
\end{itemize}

\subsection{Dla operatora: architektura}

Rysunek pokazuje przepływ. Przeglądarka nie zna klucza. Klucz do dostawcy modelu
i sekrety płatności są tylko w zmiennych środowiskowych usługi.

\begin{figure}[H]
\centering
\includegraphics[width=0.95\linewidth,height=0.42\textheight,keepaspectratio]{\dgm{p5-webapp}}
\caption{Wersja webowa: przeglądarka, BFF, \texttt{InterviewService}, dostawca
modelu i dostawca płatności. Tekst rozmowy zostaje w pamięci usługi.}
\label{fig:webapp}
\end{figure}

Elementy i ich role:

\begin{itemize}[leftmargin=*,itemsep=2pt]
  \item \textbf{\texttt{web/app}} (Next.js): strona \texttt{/interview} i ścieżka
        BFF, która przekazuje żądania do usługi bez klucza. Ciasteczko sesji
        jest po stronie BFF.
  \item \textbf{\texttt{ExitInterviewAgent.InterviewService}}: endpointy
        \texttt{/api/v1/interviews}, \texttt{/credits}, \texttt{/checkout} i
        \texttt{/webhooks/payments}. Sesja trzyma rozmowę w pamięci, jedno zadanie
        w tle na sesję. Zasada: jedna sesja otwarta na konto.
  \item \textbf{\texttt{ExitInterviewAgent.Agent}}: ta sama biblioteka co w CLI,
        z maszyną stanów i protokołem. Kafelki tworzy ten sam kod co w CLI.
  \item \textbf{authservice}: logowanie i tokeny RS256 sprawdzane po JWKS; usługa
        niczego nie podpisuje (AGENTS.md).
  \item \textbf{PostgreSQL} (\texttt{interviewdb}): identyfikatory, liczniki,
        księga kredytów, zdarzenia płatności. Schemat przez migracje, nie przez
        \texttt{EnsureCreated}.
\end{itemize}

\subsubsection{Ograniczenia architektury, które trzeba znać}

\begin{itemize}[leftmargin=*,itemsep=2pt]
  \item \textbf{Jedna instancja.} Sesje są w pamięci jednego procesu. Druga
        instancja ich nie widzi. Wdrożenie (W10) musi mieć jedną repliką albo
        przenieść sesje do wspólnego stanu.
  \item \textbf{Restart kasuje rozmowy, a zwrot kredytu nie jest zbudowany.}
        Otwarte sesje znikają, a następne żądanie dostaje 404 (ADR-0076). Plan
        (§10) obiecuje, że w tym przypadku kredyt wraca. W kodzie zwrot działa
        tylko, gdy sesja kończy się jako \texttt{failed} w czasie działania
        procesu. Nie ma uzgadniania księgi przy starcie, więc kredyt zużyty
        przed restartem \emph{nie wraca sam}. To luka do zamknięcia przed
        startem.
  \item \textbf{Wygaśnięcie bezczynności nie zwraca kredytu.} Sesja wygasła
        przez 30 minut bez żądania jest czyszczona bez przejścia przez
        \texttt{failed}, więc zwrot się nie uruchamia. To zgodne z zasadą, że
        zwrot należy się za błąd usługi, a nie za nieobecność osoby.
  \item \textbf{Brak strumieniowania.} Każda tura to jedno żądanie. Model odpowiada
        w kilka sekund; to nie zmierzono na prawdziwym modelu.
\end{itemize}

\subsection{Dla operatora: bramki kosztów}

Koszt rośnie z liczbą wywołań modelu, więc wersja webowa ma cztery hamulce
usługi i jeden hamulec po stronie dostawcy. Liczby poniżej są wartościami
domyślnymi z \repopath{appsettings.json} i są robocze, dopóki W9 nie zmierzy
kosztu rozmowy.

\begin{longtable}{@{}p{0.2\linewidth}p{0.4\linewidth}p{0.3\linewidth}@{}}
\toprule
Hamulec & Ustawienie & Co zwraca \\
\midrule
\endhead
Wyłącznik & \texttt{Interviews\_\allowbreak{}\_\allowbreak{}Enabled=false}; czytany przy każdym starcie &
503 \texttt{interviews\_\allowbreak{}disabled} \\
Limity & starty: 5 na konto i 20 na adres na godzinę; odpowiedzi: 30 na konto i
60 na adres na minutę & 429 \texttt{rate\_\allowbreak{}limited} z \texttt{Retry-After} \\
Bramka e-mail & \texttt{Interviews\_\allowbreak{}\_\allowbreak{}RequireVerifiedEmail=true} (produkcja) &
403 \texttt{email\_\allowbreak{}not\_\allowbreak{}verified} \\
Dzienny limit & \texttt{Interviews\_\allowbreak{}\_\allowbreak{}MaxStartsPerDay=20} (licznik w pamięci,
zeruje się przy restarcie) & 503 do północy UTC; \texttt{0} odrzuca wszystko \\
Limit u dostawcy & osobny workspace z miesięcznym limitem wydatków w konsoli
dostawcy & zatrzymuje rachunek; \textbf{ustawia właściciel, NOT RUN} \\
\bottomrule
\end{longtable}

Wyłącznik nie zwraca pieniędzy i nie zatrzymuje trwającej rozmowy. Limit
wydatków u dostawcy jest jedynym hamulcem rachunku. Dzienny licznik i limity
działają na żądaniach, więc nie zastąpią limitu wydatków. Pełna procedura
(wzrost kosztu, wymiana klucza, zapis incydentu) jest w runbooku
\repopath{docs/runbooks/web-app-cost-controls.md}, § 5.

Metryki nie mają etykiet (żadnego konta, sesji ani adresu). Najważniejsze:
\texttt{interviews\_\allowbreak{}started}, \texttt{interviews\_\allowbreak{}completed},
\texttt{interviews\_\allowbreak{}failed}, \texttt{interviews\_\allowbreak{}withdrawn}, \texttt{rate\_\allowbreak{}limited},
\texttt{rejected\_\allowbreak{}email\_\allowbreak{}unverified}, \texttt{tokens\_\allowbreak{}estimated} i
\texttt{model\_\allowbreak{}calls}. Średnia z histogramu to suma podzielona przez liczbę próbek.
Wczesne liczby z kilku rozmów nic nie znaczą.

\subsection{Dla operatora: kredyty i webhook}

Kredyty to księga, nie saldo w tabeli. Trzy reguły egzekwują baza i kod razem
(\texttt{CreditLedger}, ADR-0077):

\begin{enumerate}[leftmargin=*,itemsep=2pt]
  \item Zdarzenie płatności dodaje kredyty \emph{raz} (unikalny identyfikator
        zdarzenia).
  \item Sesja zużywa kredyt \emph{raz}, a zwrot dla sesji zakończonej jako
        \texttt{failed} też jest \emph{raz}.
  \item Kredyt nigdy nie schodzi poniżej zera, nawet przy równoległych żądaniach
        (na PostgreSQL każdy zapis działa w transakcji SERIALIZABLE z ponowieniem).
\end{enumerate}

Webhook płatności sprawdza podpis i tolerancję czasu (\texttt{Billing\_\allowbreak{}\_\allowbreak{}WebhookTolerance},
domyślnie 5 min). Bez ustawionego sekretu odrzuca każde żądanie. Dostawca
\texttt{fake} to dostawca dla deweloperów i testów: bez sieci i bez pieniędzy,
ale webhook przechodzi tę samą ścieżkę co produkcyjny. Dostawca \texttt{stripe}
jest napisany i testowany wyłącznie na atrapie HTTP; \textbf{test z prawdziwym
Stripe jest NOT RUN}. Karta nie trafia do usługi: płatność odbywa się na stronie
dostawcy.

\subsection{Dla operatora: co NIE zostało zrobione}

\begin{longtable}{@{}p{0.27\linewidth}p{0.68\linewidth}@{}}
\toprule
Pozycja & Stan i skutek \\
\midrule
\endhead
Przegląd prawny & NOT RUN. Lista L1--L11 w \repopath{docs/legal/CHECKLIST.md} ma
status NOT REVIEWED. Bez przeglądu nie ma płatnej rozmowy. \\
Test z prawdziwym Stripe & NOT RUN. Atrapa HTTP. Pierwsza prawdziwa płatność
jest zawsze testem na koncie testowym przed produkcyjnym. \\
Przebieg z prawdziwym modelem (W9) & NOT RUN. Brak zmierzonego kosztu, jakości
po polsku i ceny. Dopóki go nie ma, \texttt{Interviews\_\allowbreak{}\_\allowbreak{}Provider} zostaje puste
albo ustawiony na mock. \\
Limit wydatków u dostawcy & NOT RUN, decyzja właściciela. Bez niego bramki liczą
żądania, a nie pieniądze. \\
Bramka e-mail & blokuje wszystkie starty, dopóki authservice nie dodaje
\texttt{email\_\allowbreak{}verified} do tokenu (ADR-0078). \\
Wdrożenie & NOT RUN. Nic nie jest wdrożone, \texttt{flyio/web.fly.toml} istnieje,
ale nie został uruchomiony. Zakaz w AGENTS.md. \\
Zgłoszenie do sygnałów z webu & nie zbudowane (zob. wyżej). \\
Tekst zgody a kod & szkic ma trzy pola, kod jedno (zob. wyżej). \\
Polska wersja strony prywatności & nie istnieje; \texttt{/privacy} ma
tylko angielski tekst (CHECKLIST, sekcja B). \\
\bottomrule
\end{longtable}

\subsection{Dla operatora: lista kontrolna przed startem}

Każdy punkt ma sprawdzenie, które może wykonać osoba operująca usługą. Punkty
oznaczone \emph{NOT RUN} nie są wykonane w tym przewodniku.

\begin{enumerate}[leftmargin=*,itemsep=3pt]
  \item \textbf{Przegląd prawny zapisany w CHECKLIST} (nazwisko i data przy każdej
        pozycji, RELEASE-GATE pozycja 16). \emph{NOT RUN.}
  \item \textbf{Teksty polityki i regulaminu} po przeglądzie, z numerem wersji
        użytym w zgodzie (\texttt{[policy-version]} w szkicu).
  \item \textbf{Zgoda zgodna z kodem}: decyzja, czy są trzy pola, czy jedno,
        i test, że starty bez zgody nie zużywają kredytu.
  \item \textbf{Test prawdziwego Stripe} na koncie testowym, z podpisanym
        zdarzeniem i sprawdzeniem, że kredyt przychodzi raz. \emph{NOT RUN.}
  \item \textbf{Przebieg W9} na prawdziwym modelu, z kosztem, jakością i ceną
        zapisanymi w bramce wydania. \emph{NOT RUN.}
  \item \textbf{Limit wydatków} w osobnym workspace dostawcy i klucz tylko jako
        sekret wdrożenia. \emph{NOT RUN, decyzja właściciela.}
  \item \textbf{Bramka e-mail} odblokowana przez authservice (albo świadoma
        decyzja o innym mechanizmie).
  \item \textbf{Jedna replika} usługi webowej, albo przeniesienie sesji do
        wspólnego stanu.
  \item \textbf{\texttt{/health}} pokazuje integracje \texttt{interviews} i
        \texttt{interview-cost-controls} jako gotowe, a \texttt{Interviews\_\allowbreak{}\_\allowbreak{}Enabled}
        ma wartość, którą właściciel chce.
  \item \textbf{Wyłącznik sprawdzony}: ustawienie \texttt{false} daje 503, a
        istniejąca sesja dokańcza się normalnie.
\end{enumerate}

\subsection{Ćwiczenia: wersja webowa na mocku}

Ćwiczenia używają wyłącznie lokalnych komend, modelu testowego (\texttt{mock})
i atrapy płatności. Nic nie wychodzi do sieci. Wyjścia są z przebiegów w tej
sesji. Gdzie komenda się nie wykonała, jest to opisane wprost.

% ----------------------------------------------------------------------------
% Wyniki ćwiczeń wkleja się z rzeczywistego przebiegu (patrz uwaga w nagłówku).
% ----------------------------------------------------------------------------
\subsubsection{Ćwiczenie 1: cała droga webowa na testach usługi}

\textbf{Cel.} Zobaczyć, że zgoda, rozmowa, wynik, kafelki, usunięcie i wygaśnięcie
działają w usłudze, zanim ktoś otworzy przeglądarkę. Testy startują usługę w
procesie testowym: pamięć zamiast bazy, model testowy zamiast dostawcy i podpisane
tokeny testowe zamiast authservice. Nic nie wychodzi do sieci.

\textbf{Czas.} Około 35 sekund budowania i 15 sekund testów (zmierzone).

\textbf{Wymagania.} .NET SDK 10.x, pakiety NuGet przywrócone (\texttt{dotnet restore}).

\textbf{Kroki.} Z katalogu głównego repozytorium:

\begin{Verbatim}[fontsize=\scriptsize,frame=single,framesep=2mm]
dotnet test tests/ExitInterviewAgent.InterviewService.Tests \
    --no-restore --filter "FullyQualifiedName~InterviewEndpointTests" -v normal
\end{Verbatim}

\textbf{Oczekiwane wyjście} (skrócone; to są nazwy testów, które przeszły):

\begin{Verbatim}[fontsize=\scriptsize,frame=single,framesep=2mm]
  Passed ...InterviewEndpointTests.A_polish_interview_runs_from_consent_to_a_result_with_tiles
  Passed ...InterviewEndpointTests.Another_account_gets_404_for_every_route_of_a_session_it_does_not_own
  Passed ...InterviewEndpointTests.An_idle_session_expires_30_minutes_after_its_last_request_and_is_410
  Passed ...InterviewEndpointTests.The_kill_switch_refuses_new_sessions_with_503
  Passed ...InterviewEndpointTests.A_withdrawn_consent_ends_as_stopped_and_is_not_refunded
  Passed ...InterviewEndpointTests.A_provider_failure_ends_the_session_as_failed_with_503_and_refunds_the_credit
  Passed ...InterviewEndpointTests.No_interview_text_reaches_the_logs_or_the_spans
  ...
Total tests: 24
     Passed: 24
Test Run Successful.
\end{Verbatim}

\textbf{Jak sprawdzić sukces.} Ostatnia linia to \texttt{Test Run Successful}, a
\texttt{Passed: 24} zgadza się z \texttt{Total tests: 24}. Test
\texttt{A\_\allowbreak{}polish\_\allowbreak{}interview\_\allowbreak{}runs\_\allowbreak{}from\_\allowbreak{}consent\_\allowbreak{}to\_\allowbreak{}a\_\allowbreak{}result\_\allowbreak{}with\_\allowbreak{}tiles}
przechodzi pełną drogę: start po polsku, odpowiedzi, wynik z kafelkami. Test
\texttt{Another\_\allowbreak{}account\_\allowbreak{}gets\_\allowbreak{}404\_\allowbreak{}\ldots} pokazuje, że cudza sesja to 404.

\textbf{Typowe błędy.} Filtr \texttt{InterviewEndpointTests} obejmuje tylko
endpointy sesji. Testy kredytów są w ćwiczeniu 2. Bez wcześniejszego
\texttt{dotnet restore} polecenie z \texttt{--no-restore} zakończy się błędem brakujących pakietów.

\subsubsection{Ćwiczenie 2: kredyty, zakup i webhook na atrapie płatności}

\textbf{Cel.} Zobaczyć regułę „jeden kredyt na jedną rozmowę” i to, że podpisany
webhook dodaje kredyty raz, nawet przy powtórnej dostawie. Zakup idzie przez
dostawcę \texttt{fake} (bez sieci i bez pieniędzy), ale ścieżka webhooka jest ta
sama co w produkcyjnym dostawcy.

\textbf{Czas.} Około 10 sekund (zmierzone).

\textbf{Kroki.}

\begin{Verbatim}[fontsize=\scriptsize,frame=single,framesep=2mm]
dotnet test tests/ExitInterviewAgent.InterviewService.Tests \
    --no-build --filter "FullyQualifiedName~BillingEndpointTests" -v normal
dotnet test tests/ExitInterviewAgent.InterviewService.Tests \
    --no-build --filter "FullyQualifiedName~Billing"
\end{Verbatim}

\textbf{Oczekiwane wyjście} (skrócone):

\begin{Verbatim}[fontsize=\scriptsize,frame=single,framesep=2mm]
  Passed ...BillingEndpointTests.A_signed_paid_event_adds_its_credits_and_a_redelivery_adds_nothing
  Passed ...BillingEndpointTests.A_webhook_with_a_wrong_signature_is_refused_and_adds_nothing
  Passed ...BillingEndpointTests.Starting_an_interview_takes_one_credit_and_withdrawing_it_gives_nothing_back
  Passed ...BillingEndpointTests.A_session_that_fails_on_the_service_side_returns_its_credit
  Passed ...BillingEndpointTests.Checkout_with_no_provider_answers_503_billing_disabled
Total tests: 15
     Passed: 15
  ...
Passed!  - Failed:     0, Passed:    52, Skipped:     3, Total:    55
\end{Verbatim}

\textbf{Jak sprawdzić sukces.} Pierwsza komenda: 15 z 15. Druga: \texttt{Failed: 0}.
Trzy pominięte testy to \texttt{PostgresCreditTests}: wymagają działającego
PostgreSQL, którego w tej sesji nie było. Pominięcia nie są zaliczeniem. Te
trzy testy (współbieżne zużycie ostatniego kredytu i podwójna dostawa zdarzenia
na PostgreSQL) są \emph{NOT RUN}.

\textbf{Typowe błędy.} Sekret \texttt{Billing\_\allowbreak{}\_\allowbreak{}FakeWebhookSecret} nieustawiony:
dostawca \texttt{fake} losuje sekret na proces, więc zewnętrzny webhook nie
przejdzie podpisu. To zamierzone. Brak \texttt{Billing:PriceMinorUnits} daje 503
przy zakupie, nawet przy dostawcy \texttt{fake}.

\subsubsection{Ćwiczenie 3: rozmowa z modelem testowym i prośba o zakończenie}

\textbf{Cel.} Zobaczyć, jak prowadzący reaguje na krótką odpowiedź i na prośbę
„możemy już skończyć?”, na tej samej maszynie stanów co sieciowa wersja. Model
\texttt{mock} to skryptowany model testowy. Nie rozumie odpowiedzi, więc
przydziela treść do tematów mechanicznie. Ćwiczenie sprawdza przepływ, a nie
jakość rozmowy.

\textbf{Czas.} Kilka sekund.

\textbf{Wymagania.} Ćwiczenie 1 (kod zbudowany). Odpowiedzi podano potokiem,
bo nie ma terminala interaktywnego.

\textbf{Kroki.}

\begin{Verbatim}[fontsize=\scriptsize,frame=single,framesep=2mm]
printf 'tak\nWdrożenie trwało trzy tygodnie i było dobrze zorganizowane.\nsłabe\nPierwszy miesiąc: nikt nie pokazał mi procesów, sam szukałem dokumentów.\nZarządzanie: przełożony nie dawał informacji zwrotnej przez cały rok.\nNie wiem.\nMożemy już skończyć?\n' \
  | dotnet run --no-build --project src/ExitInterviewAgent.Cli -- \
      interview --provider mock --model mock --tenure 1y_3y \
      --employer acme-example --out ./iv-out --save-transcript
\end{Verbatim}

\textbf{Oczekiwane wyjście} (skrócone):

\begin{Verbatim}[fontsize=\scriptsize,frame=single,framesep=2mm]
Scripted mock model: nothing is sent anywhere. It is a test seam that understands nothing, not an interviewer.
Interviewer: Hello, and thank you for taking the time. ... I am an AI interviewer, not a person. ...
Interviewer: Let us start with your first weeks. How did you experience joining the company ...
Interviewer: Dobrze, kontynuujmy po polsku. Jak na co dzień wyglądało kierowanie pracą?
Interviewer: Jaka jedna konkretna sytuacja albo moment najlepiej to ilustruje?
Interviewer: Jakie możliwości nauki lub awansu zawodowego były tam dostępne, ...
Interviewer: Jak wynagrodzenie i świadczenia wyglądały w porównaniu z tym, ...
Interviewer: Jak można opisać sposób, w jaki ludzie tam ze sobą współpracowali?
Interviewer: Rozumiem, kończymy tutaj. Dziękuję za czas. Przygotuję zapis wyłącznie na podstawie tego, co już zostało powiedziane.

== Result ==
outcome: completed (unresponsive)
...
"piiMasked": true,
== Validation ==
valid: yes, errors: 0
submittable: yes.
Wrote: /home/user/iv-out/record.json, /home/user/iv-out/transcript.txt
...
Odrzucone: 2 (kody: pii_found)
Wrote: /home/user/iv-out/tiles/tiles.json, /home/user/iv-out/tiles/tiles.html
\end{Verbatim}

\textbf{Jak sprawdzić sukces.} Wynik to \texttt{outcome: completed (unresponsive)}.
Po odpowiedzi „Możemy już skończyć?” nie pojawia się już pytanie, tylko zamknięcie
z zapisem. To nie jest wycofanie zgody: rekord jest walidowany i zapisany. Krótka
odpowiedź „słabe” dostała jedno dopytanie („Jaka jedna konkretna sytuacja…”),
a nie zamknięcie. Walidacja daje \texttt{errors: 0}.

\textbf{Czego ten przebieg nie pokazuje.} Przypisanie tematów jest mechaniczne:
w rekordzie cytat o zarządzaniu trafił do tematu „Rozwój”, a kafelki są
zbudowane z tego, co zwrócił model testowy. Nie oceniaj jakości rozmowy po tym przebiegu.
Prawdziwa jakość po polsku to W9, którego nie wykonano (NOT RUN). Prośba o zakończenie
nie jest wycofaniem zgody; wycofanie kasuje wszystko i nie zwraca kredytu, co
sprawdza test \texttt{A\_\allowbreak{}withdrawn\_\allowbreak{}consent\_\allowbreak{}ends\_\allowbreak{}as\_\allowbreak{}stopped\_\allowbreak{}and\_\allowbreak{}is\_\allowbreak{}not\_\allowbreak{}refunded}
(ćwiczenie 1).

\textbf{Typowe błędy.} Bez \texttt{--no-build} komenda przebuduje projekt. Opcja
\texttt{--out} wskazuje katalog, nie plik: wyniki trafiły do \texttt{iv-out/record.json}
i \texttt{transcript.txt}.


}
